\documentclass{beamer}
\usepackage[utf8]{inputenc}

\usetheme{Madrid}
\usecolortheme{default}
\usepackage{amsmath,amssymb,amsfonts,amsthm}
\usepackage{txfonts}
\usepackage{tkz-euclide}
\usepackage{listings}
\usepackage{adjustbox}
\usepackage{array}
\usepackage{tabularx}
\usepackage{gvv}
\usepackage{lmodern}
\usepackage{circuitikz}
\usepackage{tikz}
\usepackage{graphicx}

\setbeamertemplate{page number in head/foot}[totalframenumber]

\usepackage{tcolorbox}
\tcbuselibrary{minted,breakable,xparse,skins}



\definecolor{bg}{gray}{0.95}
\DeclareTCBListing{mintedbox}{O{}m!O{}}{%
	breakable=true,
	listing engine=minted,
	listing only,
	minted language=#2,
	minted style=default,
	minted options={%
		linenos,
		gobble=0,
		breaklines=true,
		breakafter=,,
		fontsize=\small,
		numbersep=8pt,
		#1},
	boxsep=0pt,
	left skip=0pt,
	right skip=0pt,
	left=25pt,
	right=0pt,
	top=3pt,
	bottom=3pt,
	arc=5pt,
	leftrule=0pt,
	rightrule=0pt,
	bottomrule=2pt,
	toprule=2pt,
	colback=bg,
	colframe=orange!70,
	enhanced,
	overlay={%
		\begin{tcbclipinterior}
			\fill[orange!20!white] (frame.south west) rectangle ([xshift=20pt]frame.north west);
	\end{tcbclipinterior}},
	#3,
}
\lstset{
	language=C,
	basicstyle=\ttfamily\small,
	keywordstyle=\color{blue},
	stringstyle=\color{orange},
	commentstyle=\color{green!60!black},
	numbers=left,
	numberstyle=\tiny\color{gray},
	breaklines=true,
	showstringspaces=false,
}
%------------------------------------------------------------
%This block of code defines the information to appear in the
%Title page
\title %optional
{12.60}
%\subtitle{A short story}

\author % (optional)
{RAVULA SHASHANK REDDY - EE25BTECH11047}

 \begin{document}
	
	
	\frame{\titlepage}
	\begin{frame}{Question}
     A square matrix $\vec{B}$ is skew-symmetric if 
\begin{enumerate}[label alph*]
    \item  $\vec{B}^\top = -\vec{B}$ 
    \item  $\vec{B}^\top = \vec{B}$ 
    \item  $\vec{B}^{-1} = \vec{B}$ 
    \item  $\vec{B}^{-1} = \vec{B}^\top$ 
\end{enumerate}
\end{frame}
\begin{frame}{Solution}
     $\vec{B}$ is a skew-symmetric, then by quadratic property 
 \begin{align}
\vec{x}^\top \vec{B} \vec{x} = 0 \quad \forall \vec{x} \in \mathbb{R}^n
\end{align}

Consider two vectors $\vec{x}, \vec{y} \in \mathbb{R}^n $ :

\begin{align}
&(\vec{x}+\vec{y})^\top \vec{B} (\vec{x}+\vec{y}) = \vec{x}^\top \vec{B} \vec{x} + \vec{x}^\top \vec{B} \vec{y} + \vec{y}^\top \vec{B} \vec{x} + \vec{y}^\top \vec{B} \vec{y} \\[1mm]
&\text{Since } \vec{x}^\top \vec{B} \vec{x} = \vec{y}^\top \vec{B} \vec{y} = 0, \text{ we have:} \\[1mm]
&\vec{x}^\top \vec{B} \vec{y} + \vec{y}^\top \vec{B} \vec{x} = 0 \quad \Rightarrow \quad \vec{x}^\top \vec{B} \vec{y} = -\,\vec{y}^\top \vec{B} \vec{x} \\[1mm]
&\text{Take } \vec{x} = \vec{e_1}, \, \vec{y} = \vec{e_2} \text{ (standard basis vectors)}\\[1mm]
&\vec{e_1}^\top \vec{B} \vec{e_2} = -\, \vec{e_2}^\top \vec{B} \vec{e_1} \quad \Rightarrow \quad b_{ij} = -b_{ji} \\[1mm]
&\text{Since this holds for all } i,j, \text{ we conclude:} \\[1mm]
&\boxed{\vec{B}^\top = -\vec{B}}
\end{align}
\end{frame}
\begin{frame}{Solution}
Example: 
\[
\vec{B} = \myvec{0 & 2 & -1 \\ -2 & 0 & 3 \\ 1 & -3 & 0}
\]
\[
\vec{B}^\top = \myvec{0 & -2 & 1 \\ 2 & 0 & -3 \\ -1 & 3 & 0} = -\,\vec{B}
\]
\end{frame}
\begin{frame}[fragile]
\frametitle{C Code}
\begin{lstlisting}
    #include <stdio.h>

#define N 3  // size of the square matrix

int main() {
    int B[N][N] = {
        {0, 2, -1},
        {-2, 0, 3},
        {1, -3, 0}
    };
    
    int isSkewSymmetric = 1;  // 1 = true, 0 = false

    // Check the transpose property B^T = -B
    for (int i = 0; i < N; i++) {
        for (int j = 0; j < N; j++) {
            if (B[i][j] != -B[j][i]) {
                isSkewSymmetric = 0;
                \end{lstlisting}
                \end{frame}
\begin{frame}[fragile]
\frametitle{C Code}
\begin{lstlisting}
                break;
            }
        }
        if (!isSkewSymmetric) break;
    }
    printf("Matrix B:\n");
    for (int i = 0; i < N; i++) {
        for (int j = 0; j < N; j++) {
            printf("%4d", B[i][j]);
        }
        printf("\n");
    }
    printf("\nIs B skew-symmetric (B^T = -B)? %d\n", isSkewSymmetric);

    return 0;
}
\end{lstlisting}
\end{frame}
\begin{frame}[fragile]
\frametitle{Python Direct}
\begin{lstlisting}
    import numpy as np

# Define a skew-symmetric matrix B
B = np.array([[0, 2, -1],
              [-2, 0, 3],
              [1, -3, 0]])

# Compute the transpose
B_T = B.T

# Compute -B
neg_B = -B
\end{lstlisting}
\end{frame}
\begin{frame}[fragile]
\frametitle{Python Direct}
\begin{lstlisting}
# Check if transpose equals negative of B
is_skew_symmetric = np.allclose(B_T, neg_B)

print("Matrix B:")
print(B)
print("\nTranspose B^T:")
print(B_T)
print("\nNegative of B (-B):")
print(neg_B)
print("\nIs B skew-symmetric (B^T = -B)?", is_skew_symmetric)
\end{lstlisting}
\end{frame}

\end{document}
